% Korean edition of sections/s4_3_guidelines.tex.
\seclabel
\claim{\secnum{4.3}{sec:guidelines}\quad \code{plan()}은 목표의 섹터 한도 초과를 거부하고, \code{guidelines()}는 나머지 한도를 보고한다}

\lead{대회에는 투자 가이드라인 4개가 있고, 서버가 주문 시점에 강제하는 것은 종목 한도뿐이다(\figref{guiderules}).} 서버는 종목을 종목 한도(\code{weiLimit})보다 높게 만들 주문을 막는다. 한도는 삼성전자 40\%, 나머지 종목 15\%다. 섹터 한도는 max(2 $\times$ \code{msecs}의 시장 섹터 비중, 10\%)다. 섹터와 소형주 초과는 \code{Portfolio/Violations} 응답에 위반으로 세기만 하며, 피드의 준수 레코드(\code{cmpl})에 이번 주 회전율이 있다. 회전율은 매주 5\% 이상이어야 한다. 주간 위반은 3회까지 허용되며, 4회 이상이면 포트폴리오가 순위와 시상에서 제외된다. 한도는 대회 사이트의 가이드라인 페이지에서 가져왔다.

\lead{가이드라인에 쓰는 값은 주요 필드 4개 외의 작은 피드 필드 5개와 서버의 \code{Violations} 응답에서 나온다(\figref{guiderules}).} 스냅샷에는 \code{bs}, \code{pos}, \code{ords}, \code{msecs} 외에 \code{bm}, \code{bs_updated}, \code{cmpl}, \code{noPosPrcs}, \code{tags}가 있다. \code{state} 모듈은 이 필드들을 받은 그대로 보관하고, 델타에서 null이 아닌 값이 오면 덮어쓴다. 모든 델타에 있는 \code{bs_updated}는 \code{bs}에 병합하므로 순자산(\code{na})이 최신으로 유지된다(\ref{sec:feed}절). \code{tags}는 종목코드별로 종목에 라벨을 붙이고, \code{msecs}는 GICS 섹터마다 레코드 \code{Id}, 섹터 코드 \code{sec}, 시장 비중 \code{wei}를 가진 레코드 하나를 둔다. \code{Violations} 응답은 서버 자체의 집계를 세 부분(\code{Current}, \code{Cumulative}, \code{Details})으로 준다. 회전율 필드는 없으므로 회전율은 \code{cmpl}에서만 얻는다.

% Korean edition of figs/guiderules.tex.
\fig[wide]{guiderules}{서버는 종목 한도 초과만 막고, 나머지 가이드라인 3개는 집계만 한다}[대회 가이드라인, 각각의 강제 방식, 근거 필드, 2026년 11월 24일 09:52 KST의 값]{%
{\tabfont\begin{tabularx}{\linewidth}{>{\RaggedRight\arraybackslash}p{19mm} >{\RaggedRight\arraybackslash}p{33mm} >{\RaggedRight\arraybackslash}p{43mm} >{\RaggedRight\arraybackslash}X >{\RaggedRight\arraybackslash}p{40mm}}
\thd{규칙} & \thd{한도} & \thd{강제 방식} & \thd{09:52 값} & \thd{출처 필드}\\
종목 비중 & 종목의 \code{weiLimit}까지: 삼성전자 40\%, 나머지 15\% & 서버가 주문을 막음. 래퍼: \code{plan()}이 water-filling으로 목표마다 \code{weiLimit}, 보유하지 않은 종목은 15\%(삼성전자 40\%)로 상한. 한도를 넘는 매수는 보내기 전에 거부 & \code{StkViolCnt} 0 & \code{pos}: \code{wei}, \code{weiLimit}; \code{Violations}: \code{StkViolCnt}, \code{StockViols}\\\rowsep
섹터 비중 & max(2 $\times$ 시장 섹터 비중, 10\%)까지 & 서버가 위반을 집계. 래퍼: \code{plan()}이 \code{ValueError}로 목표 거부 & \code{SectViolCnt} 0. 최소 여유 2.01\%p(\figref{sectors}) & \code{pos}: \code{sec}, \code{wei}; \code{msecs}: \code{sec}, \code{wei}(시장 비중, \%); \code{Violations}: \code{SectViolCnt}, \code{SectorViols}\\\rowsep
소형주 & 시가총액 1조 원 미만 종목, 합계 30\%까지 & 서버가 위반을 집계. 래퍼는 \code{guidelines()}로 합계를 보고 & 1종목(\code{A090460})에 3.35\%. \code{Excess} 0 & \code{tags}(종목코드별 \code{SmallCap} 라벨); \code{Violations}: \code{SmallCapViol}(\code{TotalWei}, \code{Limit}, \code{Excess})\\\rowsep
주간 회전율 & 매주 5\% 이상. 주간 위반 4회 이상이면 포트폴리오 제외 & 서버가 피드에 주간 위반을 집계. 래퍼는 \code{guidelines()}로 보고 & \code{currWto} 4.8\%, \code{avgWto} 28.0\%, \code{wtoViols} 1 & \code{cmpl}: \code{currWto}, \code{avgWto}, \code{wtoViols}, \code{expWtoViols}\\
\end{tabularx}}}{한도는 대회 사이트의 가이드라인 페이지에서, 값은 실시간 피드(\code{pos}, \code{msecs}, \code{tags}, \code{cmpl})와 \code{Portfolio/Violations} 응답에서 가져왔다. \code{Violations}는 \code{Current}와 \code{Cumulative} 집계(09:52에 모두 0. \code{Cumulative}에는 \code{ViolDays}가 추가됨)와 \code{Details}를 돌려준다. 1주 비중은 순자산 \code{bs.na}(09:52에 837,002,427원)를 쓰며, \code{bs_updated}가 이 값을 최신으로 유지한다. \code{bm}, \code{noPosPrcs}, \code{cmpl.asr}는 보관만 하고 해석하지 않는다. \code{expWtoViols}(09:52에 1)의 뜻은 확인하지 못했다. \%p = 퍼센트포인트.}{실시간 피드와 \code{Portfolio/Violations}, DART Quant Team}


\Needspace{4\baselineskip}\lead{\code{guidelines()}는 두 출처를 응답 하나로 합치며, \code{ok} 플래그는 현재 위반 0건과 이번 주 회전율 5\%를 요구한다.} 응답에는 \code{weiLimit}까지 여유가 가장 적은 종목 5개, 보유한 모든 섹터의 비중·한도·여유(\figref{sectors}), 소형주 합계, \code{cmpl}의 회전율 값, 서버의 \code{Violations} 응답 원본이 들어 있다. \code{ok}는 서버가 센 현재 위반이 0건이고 이번 주 회전율이 5\% 이상일 때만 True다(D25).

% Korean edition of figs/sectors.tex.
\fig{sectors}{섹터 한도에 가장 가까운 것은 경기소비재로, 2.01\%p 남았다}[섹터 비중(\%)과 한도, 여유(\%p), 보유 비중이 큰 순서. 2026년 11월 24일 09:52 KST]{%
{\tabfont\begin{tabularx}{\linewidth}{>{\RaggedRight\arraybackslash}X R{19mm} R{15mm} R{14mm} R{16mm}}
\thd{GICS 섹터} & \thd{시장 비중} & \thd{포트폴리오} & \thd{한도} & \thd{여유}\\
45 IT & 52.90 & 45.74 & 105.80 & 60.06\\\rowsep
20 산업재 & 17.70 & 30.17 & 35.40 & 5.23\\\rowsep
\hl 25 경기소비재 & 5.75 & 9.49 & 11.50 & 2.01\\\rowsep
50 커뮤니케이션 서비스 & 5.00 이하 & 6.79 & 10.00 & 3.21\\\rowsep
35 헬스케어 & 5.69 & 6.54 & 11.38 & 4.84\\
\end{tabularx}}}{음영 행: 여유가 가장 적은 섹터. 한도 = max(2 $\times$ 시장 비중, 10\%). 여유 = 한도 $-$ 포트폴리오 비중. 가장 큰 섹터인 IT는 60.06\%p가 남았다. 시장 비중은 한도에서 역산했다. 커뮤니케이션 서비스는 10\% 하한이 적용되므로 시장 비중의 상한만 알 수 있다. 이 5개 섹터가 포트폴리오의 98.73\%를 차지한다.}{\code{guidelines()} 출력, DART Quant Team}


\lead{화요일 09:52 기준 이번 주 회전율은 4.8\%로 최소 5\%에 못 미쳤으므로, \code{Portfolio/Violations}가 현재 한도 위반 0건을 보고했는데도 \code{ok}는 False였다.} \code{cmpl.currWto}는 4.8\%였다. 회전율은 따로 센다. \code{cmpl.wtoViols}는 주간 회전율 위반 1회로, 규칙이 허용하는 3회 안이었다. 래퍼는 이 숫자를 보고만 하며, 무엇을 언제 거래할지는 전략에 맡긴다(N2).

\lead{시장 섹터 비중을 \code{msecs} 레코드 \code{Id}로 찾자 IT가 한도를 35.74\%p 넘은 것으로 나왔는데, 서버의 \code{SectViolCnt}는 0이었다.} 첫 버전의 \code{guidelines()}는 레코드 \code{Id}를 섹터로 썼다. 레코드 \code{Id} 1에는 섹터 10이 들어 있으므로 섹터 45를 찾으면 시장 비중이 나오지 않았고, 한도가 하한인 10\%로 떨어져 45.74\%인 IT가 한도를 넘은 것처럼 보였다. 섹터 코드 \code{sec}로 찾으면 IT 한도는 105.80\%다. 서버의 집계 0이 결함을 드러냈다.

\lead{\code{plan()}은 주문을 만들기 전에 모든 목표를 대회 규칙과 대조한다.} 순서대로 다음과 같다. 목표 합계가 \code{max_total}(100\%)을 넘으면 \code{ValueError}를 던진다. 목표 종목마다 \code{weiLimit}로 상한을 두고 초과분은 water-filling으로 다른 종목에 옮긴다. 유니버스 밖 종목이 있으면 예외를 던진다. 섹터 합계가 max(2 $\times$ 시장 비중, 10\%)를 넘으면 예외를 던진다. 마지막 검사는 \code{orders} 모듈의 \code{sector_violations()}다(D23). 그다음 진행 중인 주문의 남은 크기를 빼고 1주 비중보다 작은 주문을 뺀다(\ref{sec:oneshare}절). 매도로 매수 자금을 마련할 수 있도록 매도를 매수보다 앞에 둔다. 이 검사들은 계획 시점의 목표 비중에 적용된다. 주문이 체결되는 동안 비중을 한도 안에 붙잡아 두지는 않는다. 먼저 보낸 매도가 먼저 체결된다는 보장이 없으므로(7주가 7분 37초에 걸쳐 체결됨, \ref{sec:orders}절), 리밸런싱 중에 섹터나 합계가 잠시 한도를 넘을 수 있다(Q4, \ref{sec:validity}절). 보유하지 않은 종목은 \code{pos} 레코드가 없어 피드에 \code{weiLimit}도 없으므로, \code{plan()}은 가이드라인 값인 15\%(삼성전자 40\%)로 상한을 둔다. 주문을 보내기 전 검사도 같은 상한을 적용해 이를 넘길 매수를 거부한다. 종목별 섹터는 클라이언트가 시작할 때 한 번 불러오는 \code{Universe}(2,648종목)에서, 시장 비중은 호출할 때마다 \code{msecs}에서 가져온다.

\lead{서버는 섹터 초과를 비중이 움직인 뒤에야 집계로만 알려주므로 래퍼가 섹터를 직접 검사한다.} \code{plan()}은 목표를 거부하고 초과분을 그대로 둔다. 종목에 대해 water-filling이 하듯 섹터 초과분을 나누면 그 섹터의 모든 종목을 낮추고 다른 모든 종목을 올리게 되며, 바뀐 종목마다 하루 80건 중 1건을 쓴다. 오류는 섹터, 합계, 한도를 알려주므로 호출자가 목표를 고칠 수 있다(\figref{sectortest}).

% Korean edition of figs/sectortest.tex.
\fig{sectortest}{\code{plan()}은 산업재 +6\%p를 거부하고 +3\%p를 받아들였다}[삼성전자 비중을 HD현대중공업으로 옮겨 바꾼 산업재 비중(\%, GICS 20, 한도 35.36\%). 2026년 11월 24일 09:55 KST]{%
{\tabfont\begin{tabularx}{\linewidth}{>{\RaggedRight\arraybackslash}p{58mm} R{15mm} >{\RaggedRight\arraybackslash\leftskip=8pt}X}
\thd{목표 변경} & \thd{산업재} & \thd{\code{plan()} 결과}\\
없음(현재 비중) & 30.14 & 주문 0건\\\rowsep
\hl 산업재 +6\%p, 삼성전자 $-$6\%p & 36.14 & \code{ValueError}: \texttt{20 36.14\% > 35.36\%}. 주문 없음\\\rowsep
산업재 +3\%p, 삼성전자 $-$3\%p & 33.14 & 주문 2건: 삼성전자 3\% 매도 후 HD현대중공업 3\% 매수\\
\end{tabularx}}}{음영 행: 거부된 목표. 삼성전자는 한도가 105.80\%인 섹터 45에 속하므로 한도를 넘을 수 있는 것은 산업재뿐이었다. 아무것도 보내지 않았다.}{DART Quant Team}


\lead{산업재 한도는 09:52에 35.40\%, 09:55에 35.36\%였으므로 \code{plan()}은 호출할 때마다 모든 섹터 한도를 피드에서 읽는다.} \code{msecs}의 시장 비중은 약 12\,s마다 바뀌므로(\ref{sec:feed}절), 한 번 저장한 한도는 서버의 한도와 어긋나게 된다.

\lead{래퍼는 보유하지 않은 종목의 시가총액을 모르므로 소형주 규칙은 보고만 한다.} \code{tags}는 포트폴리오가 보유한 소형주 1종목에만 라벨을 붙였다. 따라서 목표에 새 소형주가 있어도 피드로는 알아볼 수 없고, 시장 데이터는 래퍼 밖에 둔다(N1). 30\% 한도 대비 3.35\%로 여유는 26.65\%p였다.
\verdict{\code{plan()}은 주문을 만들기 전에 합계 100\% 초과, 유니버스 밖 종목, 호출 때마다 피드에서 읽은 섹터 한도 초과를 거부한다. 이 검사는 주문이 체결되는 동안이 아니라 계획 시점의 목표 비중을 다룬다(D19, D23). \code{guidelines()}는 나머지 한도를 서버의 \code{Violations} 응답과 함께 보고하며, \code{ok}는 현재 위반 0건이고 이번 주 회전율이 5\% 이상일 때만 True다(D25).}
